\begin{figure}[H]
\centering
\includegraphics[width=1.0\linewidth]{fidelity/G7-01-ehime_trajectory_map.pdf}
\caption[冬の 1 例]{2020 年 1 月 15 日 00 UTC に南レク御荘公園から上げた場合の飛行経路。左が全体、右が着地点の拡大。和歌山の同じ日（図\ref{fig:g7-01-wakayama}）と同じく、地点のみ固定（橙の破線）と両方固定（黄）が南側に、基準（黒）と時刻のみ固定（青の破線）が北側にそろう。現行との差は \SI{11.1}{\kilo\metre}。}
\label{fig:g7-01-ehime}
\end{figure}

\begin{figure}[H]
\centering
\includegraphics[width=1.0\linewidth]{fidelity/G7-07-ehime_trajectory_map.pdf}
\caption[夏の 1 例]{同じ図の 7 月版。いったん東へ出てから西へ戻り、また東へ流れる経路で、4 通りの差は \SI{2.7}{\kilo\metre} にとどまる。拡大パネルの縮尺が冬（\SI{8}{\kilo\metre}）と夏（\SI{4}{\kilo\metre}）で違う点に注意。拡大パネルの左下に見える線は、飛行の途中で着地点の近くを一度通った部分である。}
\label{fig:g7-07-ehime}
\end{figure}
